\section{Latar Belakang}
\label{sec:latarbelakang}

Perkembangan teknologi informasi dan komunikasi telah mendorong transformasi digital di berbagai sektor, termasuk pada instansi pemerintahan. Pemerintah Kota Surabaya, sebagai salah satu kota metropolitan terdepan di Indonesia dengan predikat Smart City, terus berupaya meningkatkan kualitas pelayanan publik melalui pemanfaatan teknologi atau e-government.

Badan Pendapatan Daerah (Bapenda) Kota Surabaya merupakan instansi yang berfungsi mengelola pendapatan daerah dari sektor pajak dan mengelola keuangan daerah kota surabaya [1]. Bapenda Kota Surabaya memegang peranan vital sebagai ujung tombak dalam penghimpunan Pendapatan Asli Daerah (PAD), yang menjadi sumber utama pembiayaan pembangunan kota. Pajak Daerah merupakan salah satu sumber pendapatan asli daerah yang memiliki peranan yang sangat strategis dalam meningkatkan kemampuan keuangan daerah dan akan digunakan untuk keperluan daerah bagi sebesar-besarnya kemakmuran rakyat. Oleh karena itu, efektivitas dan efisiensi dalam pengelolaan, pelayanan, dan pengawasan pajak daerah menjadi kunci keberhasilan Bapenda.

Bapenda Kota Surabaya, sebagai instansi vital dalam penghimpunan Pendapatan Asli Daerah (PAD), sangat bergantung pada sistem informasi untuk mengelola data perpajakan yang kompleks. Namun, sistem monitoring yang digunakan sebelumnya menghadapi tantangan performa yang signifikan. Aplikasi tersebut menjadi sangat lambat karena database yang kurang terorganisir serta pendekatan load data yang tidak efisien, di mana keseluruhan data dimuat secara serentak saat aplikasi dibuka. Hal ini menghambat efektivitas kerja pegawai dalam melakukan analisis dan pengawasan.

Menyadari kendala tersebut, diperlukan adanya sebuah pembaruan sistem yang tidak hanya menata ulang database, tetapi juga mengimplementasikan arsitektur aplikasi yang modern dan lebih baik secara performa. Oleh karena itu, pengembangan sistem Monitoring Pendapatan Daerah (MonPD) yang baru diinisiasi sebagai solusi. Proyek ini berfokus pada migrasi ke database yang lebih baik serta penerapan teknik load data secara asinkron, sehingga aplikasi hanya akan mengambil data yang dibutuhkan pada saat itu saja. Sebagai mahasiswa Departemen Teknik Komputer, kami melihat peluang besar ini untuk menerapkan ilmu di bidang teknologi informasi dalam membangun sistem informasi yang cepat, efisien, dan andal untuk mendukung program digitalisasi di Bapenda Kota Surabaya.
